\subsection{Dimensionierung einer Komparatorschaltung mit Diode}
% Alt: Aufgabe 7
\Aufgabe{
    \label{Aufg7}
    Dimensionieren Sie die Beschaltung des Komparators so, dass das unten skizzierte Übertragungsverhalten (grün: $U_\mathrm{E}$, blau: $U_\mathrm{A}$) entsteht.
    Die Durchlassspannung der Diode wird mit $U_\mathrm{D} = \mathrm{0,7\,V}$ angenommen.
    Die maximale Verlustleistung des Widerstandes beträgt 1/4\,W.
    \begin{figure}[H]
        \centering
        % Linkes Bild
        \begin{subfigure}[b]{0.45\textwidth}
            \centering
            \input{Tikz/src/FigKomparatorDiode.tex}\alttext{Das Bild zeigt ein elektrisches Schaltbild. Links ist eine Spannungsquelle mit der Bezeichnung \(U_\mathrm{e}\), die nach unten zur Masse zeigt. Von dieser Quelle geht eine Leitung zu einem Operationsverstärker mit der Bezeichnung N1, dessen Ausgang nach rechts führt. Vom Ausgang des Operationsverstärkers geht eine Leitung zu einem Widerstand mit der Beschriftung \(R\). Der andere Anschluss des Widerstands ist mit zwei in Reihe geschalteten Dioden verbunden, die nach unten zur Masse führen. Die erste Diode ist mit \(D_1\) beschriftet, die zweite Diode darunter mit \(D_2\). Rechts neben den Dioden ist die Ausgangsspannung \(U_\mathrm{a}\) angegeben, ebenfalls mit einem Pfeil nach unten zur Masse.}
        \end{subfigure}
        % Rechtes Bild
        \begin{subfigure}[b]{0.5\textwidth}
            \centering
            \includesvg[width=0.95\textwidth]{Bilder/Aufgaben/FigKomparatorDiode}\alttext{Rechts neben der Schaltung ist ein Diagramm dargestellt, das das gewünschte Ausgangsverhalten zeigt. Auf der x-Achse ist die Zeit \(t\) in Sekunden dargestellt, auf der y-Achse die Spannung \(U\) in Volt. Die grüne Linie stellt die Eingangsspannung \(U_\mathrm{E}\) dar, die sich als periodisches Signal mit ansteigenden und abfallenden Spannungswerten verhält. Die blaue Linie stellt die Ausgangsspannung \(U_\mathrm{A}\) dar, die sich wie ein periodisches Rechtecksignal mit konstanten Werten verhält.}
        \end{subfigure}
        \label{fig:FigKomparatorDiode}
    \end{figure}
}


\Loesung{
    Formel:
    \begin{align*}
        P_\mathrm{R} = \frac{U_\mathrm{R}^2}{R}
    \end{align*}


    Spannung am Widerstand $R$:
    \begin{align*}
        U_\mathrm{R} & = U_\mathrm{a} - U_\mathrm{D}      \\
                     & = 14\,\mathrm{V} - 0,7\,\mathrm{V} \\
                     & = 13,3\,\mathrm{V}
    \end{align*}

    Berechnung des Widerstands $R$:
    \begin{align*}
        P_\mathrm{R} & = \frac{U_\mathrm{R}^2}{R} \leq \frac{1}{4}\,\mathrm{W} \\
        R            & = \frac{U_R^2}{P_{\text{max}}}
        = \frac{(13,3\,\mathrm{V})^2}{0,25\,\mathrm{W}}
        = \frac{176,89\,\mathrm{V}^2}{0,25\,\mathrm{W}}
        = 707,56\,\Omega
    \end{align*}

    Normierter Widerstandswert:
    \begin{align*}
        R & = 710\,\Omega
    \end{align*}

    Überprüfung der Verlustleistung:
    \begin{align*}
        P_\mathrm{R} & = \frac{(13,3\,\mathrm{V})^2}{710\,\Omega} \approx 0,25\,\mathrm{W}
    \end{align*}

    Der Widerstand $R$ sollte den Wert:
    \begin{align*}
        R & = 710\,\Omega
    \end{align*}
    haben, um die maximale Verlustleistung nicht zu überschreiten und das gewünschte Übertragungsverhalten zu gewährleisten.

    Siehe: Abschnitt \ref{fig: Verschaltung Verstaerker} und Tabelle \ref{tab:Grundschaltungen}

}